% Part: computability
% Chapter: computability-theory
% Section: universal-part-function

\documentclass[../../../include/open-logic-section]{subfiles}

\begin{document}

\olfileid{cmp}{thy}{uni}
\olsection{ہمہ گیر جزوی قابلِ حساب فنکشن}

\begin{thm}
\ollabel{thm:univ-comp}
اک ہمہ گیر جزوی قابلِ حساب فنکشن~$\fn{Un}(e,x)$ موجود اے۔
ہور لفظاں وچ، اک فنکشن $\fn{Un}(e,x)$ موجود اے جیہدے لئی:
\begin{enumerate}
\item $\fn{Un}(e,x)$ جزوی قابلِ حساب اے۔
\item جے $f(x)$ کوئی وی جزوی قابلِ حساب فنکشن اے، تاں
اک قدرتی عدد $e$ موجود اے ایداں کہ $f(x) \simeq \fn{Un}(e,x)$
ہر~$x$ لئی۔
\end{enumerate}
\end{thm}

\begin{proof}
$\fn{Un}(e,x) \simeq U(\umin{s}{T(e,x,s)})$ رکھو، جتھے $U$ تے
$T$ کلینی دے نارمل روپ دے قضیے (\olref[nfm]{thm:normal-form})
والے نیں۔
\end{proof}

\begin{explain}
ایہ صرف ایس گل دا ٹھیک ٹھیک بیان اے کہ ساڈے کول جزوی
قابلِ حساب فنکشناں دی اک موثر گنتی وار فہرست اے۔ خیال
ایہ اے کہ جے اسی $f_e$ اوس فنکشن لئی لکھئیے جیہدی تعریف
$f_e(x) = \fn{Un}(e,x)$ اے، تاں سلسلے $f_0$، $f_1$، $f_2$،
\dots وچ سارے جزوی قابلِ حساب فنکشن شامل نیں، تے خاصیت
ایہ اے کہ $f_e(x)$ نوں $e$ تے~$x$ دے لحاظ نال ``یکساں''
طریقے اُتے حسابیا جا سکدا اے۔ سادگی لئی اسی دو آرگومنٹاں
والا فنکشن ورت رہے آں جیہڑا اک آرگومنٹ والے فنکشناں لئی
ہمہ گیر اے، پر عدداں دے سلسلیاں دی کوڈنگ نال ایہنوں
ہور ودھ آرگومنٹاں تک آسانی نال عام کیتا جا سکدا اے۔
% PNB-CMP-012: اصل final g نوں recursive کہندا؛ partial f توں صرف partial g لازم،total نہیں۔
مثلاً، جے $f(x,y,z)$ اک $3$-آرگومنٹ والا جزوی بازگشتی
فنکشن اے، تاں فنکشن $g(x) \simeq f((x)_0, (x)_1, (x)_2)$
اک آرگومنٹ والا جزوی بازگشتی فنکشن اے۔
\end{explain}

\end{document}
